\section{Support-orbit rigidity and selected ranks}\label{sec:orbits}
% Reusable proof text, independently derived. Suggested bibliography keys:
% Nomura1994 (correct DOI 10.1007/BF02108300), Hatcher2002,
% Baez2002 (DOI 10.1090/S0273-0979-01-00934-X), FK1994.
% The real affine-kernel lemma below replaces the longer determinantal proof.

\subsection{Saturation and support orbits}
For a simple Euclidean Jordan algebra $V$ of rank $r$ and Peirce
constant $a$, write $\mathcal I_p(V)$ for its manifold of rank-$p$
idempotents. Its dimension is $ap(r-p)$ and its tangent space at an idempotent $c$
is $V(c,1/2)$. These are classical facts about idempotent manifolds
\cite{FK1994,Nomura1994}. We use them to identify the additional restriction
imposed by globally smooth factor maps.

\begin{theorem}[Global saturation]\label{thm:global-saturation}
Let $C\subset\R^{n+1}$ be a compact convex body, $n\geq1$, containing the origin in its interior,
and suppose $P=\partial C$ and $D=\partial C^\circ$ are strictly convex
$C^1$ hypersurfaces. Suppose the full slack has a finite globally labelled
factorization
\[
 1-\ip{x}{z}=\sum_i\ip{X_i(x)}{Y_i(z)},\qquad (x,z)\in P\times D,
\]
where $X_i:P\to K_i$ and $Y_i:D\to K_i$ are $C^1$ and $K_i$ is a simple
symmetric cone with rank $r_i$ and Peirce constant $a_i$. Put
$c_i=a_i\lfloor r_i^2/4\rfloor$, with $c_i=0$ for rays.
If $\sum_i c_i=n$, each positive-capacity factor has constant balanced
complementary contact ranks $p_i,r_i-p_i$, and the joint support map
\[
 P\longrightarrow\prod_{i:c_i>0}\mathcal I_{p_i}(V_i),
 \qquad x\longmapsto(\operatorname{supp}X_i(x))_i
\]
is a finite covering. There is exactly one positive-capacity factor.
It is the Lorentz cone $Q_{n+2}$, except that $n=2$ also permits
$\mathbb S_+^3$ at the level of the support-orbit topology.
\end{theorem}
\begin{proof}
Let $z$ be the unique polar contact of $x$. The contact correspondence is
a homeomorphism; no differentiability of that correspondence is needed.
The tangent spaces are $T_xP=z^\perp$ and $T_zD=x^\perp$.
Their Euclidean pairing is nondegenerate: its left kernel consists of
vectors in $z^\perp\cap\operatorname{span}\{x\}$, which is zero because
$\ip{x}{z}=1$. Mixed differentiation in independent boundary charts gives
\[
 \ip{u}{v}=\sum_i H_i(u,v),\qquad
 H_i(u,v)=-\ip{dX_i(x)[u]}{dY_i(z)[v]}.
\]
The cross-Peirce argument of Lemma~\ref{lem:peirce-channel}, which for its
rank bound uses only the two independent derivative images, gives
\[
 \operatorname{rank}H_i\leq a_ip_iq_i\leq c_i,
 \quad p_i=\jrank X_i(x),\quad q_i=\jrank Y_i(z).
\]
Nondegeneracy and $\sum_i c_i=n$ force equality throughout. For every
positive-capacity factor, this gives
$p_i+q_i=r_i$, the ranks are balanced, and $\operatorname{rank}H_i=c_i$.
Ranks are lower semicontinuous along the contact graph. Since the two
ranks sum to $r_i$, each is locally constant, hence constant by
connectedness.

In the following support-differential argument all indices satisfy
$c_i>0$; zero-capacity factors have $H_i=0$ and play no role.
Let $e_i(x)=\operatorname{supp}X_i(x)$. This is $C^1$: locally the zero
spectral cluster is separated from the positive spectrum, and a smooth
scalar cutoff equal to zero on the former and one on the latter gives
$e_i$ by spectral functional calculus. Differentiating
$e_i\circ e_i=e_i$ and $X_i\circ e_i=X_i$ shows that
\[
 de_i[u]\in V_i(e_i,1/2),\qquad
 L_{X_i}de_i[u]=\tfrac12\operatorname{proj}_{V_i(e_i,1/2)}dX_i[u].
\]
Here $L_{X_i}$ is invertible on the half-Peirce space: in a frame adapted
to the positive support it acts on a cross channel by $\lambda_j/2>0$.
Consequently $de_i[u]=0$ precisely when the cross-Peirce component of
$dX_i[u]$ vanishes. In particular
\[
 c_i=\operatorname{rank}H_i\leq\operatorname{rank}de_i\leq c_i.
\]
If every $de_i$ kills $u$, every $H_i(u,v)$ is zero, so $u=0$.
The joint support differential is therefore injective, and its domain and
codomain have the same dimension $n$. The joint map is a local
$C^1$ diffeomorphism. Its image is open and, by compactness, closed in the
connected target. It is surjective and proper, hence a finite covering.

We give the topological classification explicitly. For $n\geq2$, the
source $P\cong S^n$ is simply connected, so this is the universal cover of
the target. There can be no circle factor, since it would make the target
fundamental group infinite. The remaining factors have compact simply
connected universal covers. The universal cover of their product is the
product of those covers. If there were at least two positive-dimensional
factors, the mod-two top cohomology class of one factor, tensored with the
unit classes of the others, would give nonzero cohomology in a degree
strictly between zero and $n$. The K\"unneth formula contradicts the
cohomology of $S^n$. Thus there is one factor. For $n=1$, the same conclusion
follows directly from the sum of the positive integer dimensions.

The simple-algebra classification now leaves the following cases.
A spin factor has primitive-idempotent orbit $S^{m-2}$ and therefore gives
$Q_{n+2}$. For $H_r(\R)$ the orbit is the real Grassmannian
$\operatorname{Gr}_p(\R^{p+q})$, $p+q=r$. At $r=3$ its balanced orbit is
$\mathbb{RP}^2$, covered by $S^2$. At $r\geq4$, balanced ranks satisfy
$p,q\geq2$. The oriented Grassmannian is simply connected by the exact
sequence of
\[
 SO(p)\times SO(q)\longrightarrow SO(p+q)
 \longrightarrow\operatorname{Gr}_p^+(\R^{p+q}):
\]
the induced map on fundamental groups onto $\pi_1(SO(p+q))=\mathbb Z_2$
is surjective. Its kernel is nonzero: use the diagonal element when
$p,q\geq3$, and an even nonzero winding in an $SO(2)$ factor otherwise.
Exactness therefore gives nonzero $\pi_2$ of the oriented Grassmannian.
Its dimension $pq\geq4$ precludes a sphere.

The complex Grassmannian is simply connected by the fibration with
isotropy $U(p)\times U(q)$ in $U(p+q)$. The map on $\pi_1$ is addition
$\mathbb Z^2\to\mathbb Z$, so its kernel gives nonzero $\pi_2$ of the
Grassmannian. For $r\geq3$ its dimension $2pq\geq4$ precludes a sphere.
For quaternionic Grassmannians, the groups $Sp(k)$ are simply connected
and $\pi_3(Sp(k))=\mathbb Z$, with standard inclusions preserving a
generator. These facts follow inductively from
$Sp(k-1)\to Sp(k)\to S^{4k-1}$ and $Sp(1)=S^3$.
The map on $\pi_3$ of the isotropy inclusion is again addition
$\mathbb Z^2\to\mathbb Z$, giving nonzero $\pi_4$ of the Grassmannian.
Its dimension $4pq\geq8$ precludes a sphere. The classical-group and
Grassmannian fibrations used here are described in
\cite[Examples~4.53--4.55]{Hatcher2002}.
Finally the balanced idempotent orbit of $H_3(\OO)$ is the Cayley plane:
rank-one idempotents are its points and complementation identifies the
rank-two orbit with it~\cite[Section~3.4]{Baez2002}.
Its cells have dimensions $0,8,16$, so $H^8(-;\mathbb Z)=\mathbb Z$;
see~\cite[Example~4.47]{Hatcher2002}. It is not a sphere.
Order-two real, complex, and quaternionic Hermitian algebras are already
spin factors. This exhausts the classification.
\end{proof}

\subsection{Eliminating the real order-three exception}
The apparent exceptional possibility is ruled out by the full ball slack.
The following elementary observation avoids a classification of linear
determinantal representations.

\begin{example}[Local constant rank need not extend globally]
\label{ex:local-real-three-pencil}
The affine pencil
\[
 L(x)=\begin{pmatrix}
 1+x_3&x_1&x_2\\x_1&1-x_3&0\\x_2&0&1-x_3
 \end{pmatrix},\qquad
 \det L(x)=(1-x_3)(1-\norm{x}^2),
\]
is positive semidefinite of rank two on $S^2\setminus\{(0,0,1)\}$.
Indeed, its lower $2\times2$ block is positive definite and its Schur
complement is zero there. Its kernel line is
\[
 \operatorname{span}\left(1,-\frac{x_1}{1-x_3},
                              -\frac{x_2}{1-x_3}\right),
\]
so its kernel map is a local diffeomorphism onto an affine chart of
$\mathbb{RP}^2$ by stereographic coordinates. At the north pole,
$L(0,0,1)=\operatorname{diag}(2,0,0)$ has rank one.
Thus affine behavior and constant rank on an open sphere patch do not
justify global constant rank or a common kernel. This example only
invalidates that local-to-global shortcut. Proposition~\ref{prop:no-real-three-saturation}
below excludes the global slack factorization using the full $C^1$
factor maps and finite affine gluing.
\end{example}

\begin{lemma}[Injective kernels of an affine sphere pencil]
\label{lem:affine-kernel-injective}
Let $s\geq2$ and $X(x)=C_0+\sum_{j=1}^s x_jC_j$ be a real symmetric
matrix pencil, positive semidefinite of nullity one for every
$x\in S^{s-1}$. If its kernels do not share a nonzero vector, its
kernel-line map is injective.
\end{lemma}
\begin{proof}
Suppose distinct $x,y$ have the same kernel line, containing $k\ne0$.
The affine scalar function $f(v)=k^TX(v)k=a+b\cdot v$ is nonnegative on
the unit sphere and vanishes at $x$ and $y$. Nonnegativity says
$a\geq\|b\|$. If $b\ne0$, a zero is possible only when
$a=\|b\|$ and $v=-b/\|b\|$, giving only one zero. Thus $a=b=0$.
For every $v$, positivity and $k^TX(v)k=0$ imply $X(v)k=0$, contradicting
the absence of a common nonzero kernel vector.
\end{proof}

\begin{lemma}[Finite affine selection on a sphere]
\label{lem:finite-affine-selection}
Let $n\geq2$ and let $F:S^n\to W$ be $C^1$, with $W$ finite dimensional.
Suppose a dense open subset is the union of finitely many nonempty open
sets $U_i$, and $F|_{U_i}=L_i|_{U_i}$ for ambient affine maps
$L_i:\R^{n+1}\to W$. Then $F$ is the restriction of one affine map.
\end{lemma}
\begin{proof}
Merge sets with identical $L_i$ and put $E_i=\overline{U_i}$. Continuity
of values and tangent derivatives gives $j^1F=j^1(L_i|_{S^n})$ on $E_i$.
If distinct affine maps have the same tangent first jet at $z\in S^n$,
their difference is $M(z\cdot x-1)$ for some nonzero $M\in W$.
Such a difference has matching first jet at no other sphere point:
matching derivatives force the other normal line to be $\R z$, and
matching values exclude $-z$. Thus each intersection $E_i\cap E_j$ lies
in at most one point. Remove the resulting finite set $Z$. The connected
space $S^n\setminus Z$ is partitioned into the disjoint relatively closed
sets $E_i\setminus Z$; each is also relatively open because the partition
is finite. Only one is nonempty. But the closure of a nonempty open $U_i$
cannot be finite, so there was only one distinct affine map. Density
finishes the proof.
\end{proof}

\begin{proposition}[No real order-three saturation for a ball]
\label{prop:no-real-three-saturation}
There is no factorization
\[
 1-x\cdot y=\operatorname{tr}(X(x)Y(y))+\sum_{j=1}^q a_j(x)b_j(y),
 \qquad x,y\in S^2,
\]
with $X,Y:S^2\to\mathbb S_+^3$ of class $C^1$, continuous nonnegative
$a_j,b_j$, finite $q$, constant complementary ranks $2,1$, and a
rank-one support map that is a local diffeomorphism onto $\mathbb{RP}^2$.
\end{proposition}
\begin{proof}
Exchange $X,Y$ if necessary so that $Y$ has rank one. Every diagonal
summand is zero, hence $a_j(x)b_j(x)=0$. For each of the finitely many
continuous nonnegative functions $f=a_j,b_j$, the set
$\{f>0\}\cup\operatorname{int}\{f=0\}$ is dense open. Intersect these
sets and partition the result by their finitely many positivity/zero
patterns. On each nonempty open pattern set $U$, for each $j$ at least
one of $a_j|_U,b_j|_U$ is identically zero. Therefore all ray terms vanish
for every pair $(x,y)\in U\times U$.

The rank-one support map of $Y$ sends $U$ onto a projective open set.
Consequently the matrices $Y(y)$, $y\in U$, span $\mathbb S^3$: an
annihilating symmetric matrix would have a quadratic form zero on an open
set of projective directions, and hence zero identically. Choose six
$y_\ell\in U$ whose matrices are a basis. The equations
\[
 \operatorname{tr}(X(x)Y(y_\ell))=1-x\cdot y_\ell,
 \qquad x\in U,
\]
force $X|_U$ to be an ambient affine pencil. There are only finitely many
pattern sets, so Lemma~\ref{lem:finite-affine-selection} makes $X$ one
affine pencil on the whole sphere.

Complementarity identifies the kernel line of $X(x)$ with the support
line of $Y(x)$. This map is a finite covering $S^2\to\mathbb{RP}^2$,
hence the universal double cover. Its image contains all lines, so its
kernels have no common nonzero vector. But
Lemma~\ref{lem:affine-kernel-injective} makes the map injective,
contradicting its two-point fibers.
\end{proof}

\begin{corollary}[Ball saturation classification]
\label{cor:ball-global-saturation}
A finite globally bi-$C^1$ symmetric-cone factorization of the full slack
of $\B_2^s$, $s\geq2$, satisfying
$\sum_i a_i\lfloor r_i^2/4\rfloor=s-1$ has exactly one
positive-capacity factor, namely $Q_{s+1}$. Remaining factors, if any,
are rays.
\end{corollary}
\begin{proof}
Apply Theorem~\ref{thm:global-saturation} and exclude its sole additional
case by Proposition~\ref{prop:no-real-three-saturation}.
\end{proof}


\subsection{The selected-rank frontier for a fixed dictionary}
For the dictionary of Section~\ref{sec:rank}, define
\begin{equation}\label{eq:smooth-kappa}
 \kappa_{\mathscr K}(s)=
 \begin{cases}
 1,&s=B+1\text{ and the dictionary contains }Q_{B+2},\\
 \lceil s/B\rceil,&\text{otherwise}.
 \end{cases}
\end{equation}
Here ``contains'' allows linear cone isomorphism. The formula also gives
one when $s\leq B$. A dictionary of Hermitian matrices over a fixed field
$\mathbb F$ with order cap $R$ has $B=\delta(R-1)$ and the exceptional
branch occurs exactly when $R=2$ and $s=\delta+1$.

\begin{theorem}[Exact selected-rank frontier]\label{thm:smooth-rank}
For globally bi-$C^1$ full-slack factorizations of $\B_2^s$ through a
fixed finite repeatable symmetric-cone dictionary, the least possible
value of $\sup_v\sum_i\jrank Y_i(v)$ is $\kappa_{\mathscr K}(s)$.
The upper bound is attained by a relative-Slater affine lift with globally
smooth boundary selections and genuine whole-slice certificates.
\end{theorem}
\begin{proof}
At every contact, with primal and dual ranks $p_i,q_i$,
\[
 s-1\leq\sum_i a_ip_iq_i
 \leq\sum_i a_i(r_i-q_i)q_i
 \leq B\sum_iq_i.
\]
Rays supply zero in the first two sums and cannot make any equality
saturated. Unless $s-1=qB$ for an integer $q\geq1$, the lower bound
already rounds to $\lceil s/B\rceil$. In the remaining case suppose
$Q(v)=\sum_iq_i(v)\leq q$ everywhere. Equality forces $Q\equiv q$;
lower semicontinuity makes each $q_i$ constant. Each nonzero dual factor
has $q_i=1$, $a_i(r_i-1)=B$, and $p_i=r_i-1$. Its primitive support
orbit has dimension $B$. The support-differential argument above gives
a finite covering from $\Sph^{qB}$ to the product of these $q$
primitive-idempotent orbits.

If $q>1$, a circle target gives infinite fundamental group. Otherwise
the compact universal covers of the positive-dimensional target factors
have nonzero intermediate product cohomology, ruling out a sphere
source. If $q=1$, the simple-algebra classification identifies the
primitive orbit as a sphere for spin factors, real projective space for
real Hermitian factors, complex or quaternionic projective space in the
other classical cases, and the Cayley plane in the exceptional case.
Complex and quaternionic projective spaces have intermediate cohomology
unless their matrix order is two; the Cayley plane has its nonzero
middle cohomology. These order-two cases are spin factors. For real
order $R\geq3$ the remaining map is the double covering
$\Sph^{R-1}\to\mathbb{RP}^{R-1}$. All other dual factors vanish
identically. Pairing the primal sheet with a basis chosen from the
rank-one dual sheet makes it an ambient affine pencil, since rank-one
matrices with directions in a projective open set span the symmetric
matrices. Its kernel-line map has no common nonzero kernel and hence
is injective by Lemma~\ref{lem:affine-kernel-injective}, contradicting
the double cover. Thus the only saturated possibility is the stated
matching spin factor. In every other case some rank is at least $q+1$.

For attainment outside the exceptional branch, choose a dictionary member
attaining $B$ and partition all $s$ coordinates into
$g=\lceil s/B\rceil$ nonempty groups embedded in its primitive half-Peirce
space. With $d=e-c$ as in Lemma~\ref{lem:perspective}, impose
\begin{equation}\label{eq:fixed-tail}
 X_G=t_Gc+\sqrt2 w_G+d\succeq0,
 \qquad\sum_Gt_G=1.
\end{equation}
These constraints say $t_G\geq\norm{w_G}^2$, so the projection is the
ball and the slice is strictly feasible. Its polynomial boundary sheet
has $t_G=\norm{w_G}^2$. At support $v=(g_G)_G$ the rank-one completion
\[
 Y_G=\tfrac12 c-\frac{g_G}{\sqrt2}+P(g_G)c
     =\tfrac12 P(c-\sqrt2 g_G)c
\]
has pairing $t_G/2-\ip{w_G}{g_G}+\norm{g_G}^2/2$.
It is nonzero of rank one even when $g_G=0$. Summing proves the genuine
whole-slice slack identity and gives selected rank $g$ everywhere.
For the exceptional branch use the direct Lorentz trace section, whose
boundary primal and dual factors have rank one. This proves attainment.
\end{proof}

\begin{theorem}[Sequential smooth ranks]\label{thm:smooth-product-rank}
Suppose the full rows of $\prod_a\B_2^{s_a}$ have globally bi-$C^1$
factors through $\mathscr K$. There is a simultaneous contact $u$ such
that, for every choice of positive weights,
\begin{equation}\label{eq:smooth-product-rank}
 \sum_i\jrank\left(\sum_a\lambda_aY_i^a(u_a)\right)
 \geq\sum_a\kappa_{\mathscr K}(s_a).
\end{equation}
The same bound holds for total primal nullity there. These bounds are
sharp over globally selected factorizations, and their upper constructions
are actual affine lifts with genuine certificates.
\end{theorem}
\begin{proof}
Use the successive joins $c_i$ from Theorem~\ref{thm:product-existence}.
After fixing earlier primal coordinates at contact, cylindrical
complementarity puts the primal factors in $V_i(e_i-c_i,1)_+$.
Fix future primal coordinates arbitrarily and apply the next one-ball
row to the compressed dual $P(e_i-c_i)Y_i^a$. The compressed maps remain
globally $C^1$ and satisfy the exact full two-point row identity.
Lemma~\ref{lem:join} identifies their total rank with the next support
increment.

To apply Theorem~\ref{thm:smooth-rank}, note that proper faces have
strictly smaller per-rank capacity than their parent simple non-ray
cone: classical faces have smaller order, spin faces are rays, and
Albert proper non-ray faces have capacity eight rather than sixteen.
Thus each remaining face has capacity at most $B$. If equality is
needed in the rank chain of that theorem, the face must be an unreduced
original cone attaining $B$. In particular, a newly created proper face
cannot introduce a matching spin exception at capacity $B$. The same
proof therefore gives at least $\kappa_{\mathscr K}(s_a)$ new rank at
some next contact. Iteration and Lemma~\ref{lem:join} prove the bound;
complementarity gives nullity. Separate copies of~\eqref{eq:fixed-tail}
and the exceptional direct sections give equality at every contact.
\end{proof}

In particular, under an integer real-dimension cap $d\geq3$ on arbitrary symmetric
factors, the exact selected aggregate-rank frontier is
\[
 \sum_a h_d(s_a),\qquad
 h_d(s)=\begin{cases}1,&d\geq s+1,\\
 \lceil s/(d-2)\rceil,&d<s+1.
 \end{cases}
\]
Indeed the per-rank capacity is at most $d-2$, with equality only for a
spin factor of dimension $d$. The equality analysis in
Theorem~\ref{thm:smooth-rank} gives the lower bound, and separate grouped
Lorentz sections attain it. This concerns the chosen global sheets.
The nonsmooth norm-tree certificates in Example~\ref{ex:tree-regularity}
show why the same smooth premium cannot be imposed on arbitrary
affine-lift certificate fibers.

There is also a factor-aware refinement. At stage $a$ of sequential
compression, suppose the non-ray faces have ranks $R_{ia}$ and Peirce
constants $\alpha_{ia}$. Define
\[
 g_a=\min\left\{\sum_iq_i:
 q_i\in\{0,\ldots,R_{ia}\},\quad
 \sum_i\alpha_{ia}q_i(R_{ia}-q_i)\geq s_a-1\right\}.
\]
The mixed-channel inequality gives at least $g_a$ new non-ray rank at
every next contact. Hence one can replace the summand in
\eqref{eq:smooth-product-rank} by
$\max\{g_a,\kappa_{\mathscr K}(s_a)\}$. This is an adaptive lower bound,
not an asserted exact frontier. For at most $L$ faces of a common Peirce
constant $\alpha$ and order at most $R$, a fixed total dual rank
$Q=uL+t$, $0\leq t<L$, has capacity at most
\[
 \alpha\big[(L-t)u(R-u)+t(u+1)(R-u-1)\big].
\]
To prove it, pad with zero ranks and increase every order to $R$.
If two ranks differ by at least two, transferring one unit from the
larger to the smaller increases their sum of $q(R-q)$. Iteration
balances all ranks to $u$ or $u+1$, proving the formula. Inverting this
finite envelope gives a convenient explicit lower bound on $g_a$.

\begin{corollary}[Dimension-capped smooth symmetric-cone resources]
\label{cor:smooth-symmetric-resources}
Consider globally bi-$C^1$ full-slack factorizations of $\B_2^s$ over
simple symmetric cones of real dimension at most an integer $d\geq3$, with optional
rays. Let $L$ count non-ray factors, $D$ the entire ambient dimension,
and $\rho$ the entire ambient Jordan rank. If $d<s+1$, their simultaneous
exact minima are
\[
 L_{\min}=\left\lceil\frac{s}{d-2}\right\rceil,\qquad
 D_{\min}=s+2L_{\min},\qquad\rho_{\min}=2L_{\min}.
\]
If $d\geq s+1$, the exact triple is $(1,s+1,2)$.
The standard ambient product barrier has parameter $\rho$.
\end{corollary}
\begin{proof}
For every simple non-ray cone,
$a\lfloor r^2/4\rfloor\leq m-2$, with equality only for rank-two
spin factors; substitution of $m=r+ar(r-1)/2$ verifies strictness
for $r\geq3$. Write $T=\sum_{i\text{ non-ray}}(m_i-2)$.
Local curvature gives $T\geq s-1$. If equality holds, all factors are
spin and Corollary~\ref{cor:ball-global-saturation} forces the single
factor $Q_{s+1}$. Thus in the strict-cap case $T\geq s$.
Since $T\leq(d-2)L$, $D\geq T+2L$, and $\rho\geq2L$, the displayed
lower bounds follow. The grouped Lorentz construction attains them
simultaneously. In the other case local curvature gives $T\geq s-1$,
$L\geq1$, and hence $D\geq s+1$, $\rho\geq2$; the direct Lorentz
section attains equality.
\end{proof}
